\MajorSection{القسم الثامن}

\Couplet{١٨٦}{لا جنة ولا جحيم؛ إنهما أحلام العقول الطفولية؛}{وأدوات كاهن الأوثان الماكر، يخوّف بها الحمقى الذين يعميهم دهاؤه.}{}

\Couplet{١٨٧}{تعلّم من أرواح القدماء العظيمة أن تضع قدمك على الجنة والجحيم؛}{وأن تجد جحيمك وجنتك في الحياة، بحسب ما تسيء استخدامها أو تحسن.}{}

\Couplet{١٨٨}{هكذا رأى اليهودي الجريء الذي تجرأ، بالصمت المدروس، أن يطرح أرضًا}{أوركوس وهاديس، بلاد الظلال، الليل الكئيب للنهار الإنساني.}{}

\Couplet{١٨٩}{الموت النهائي صعب على القلب؛ والكائن يود ألا ينتهي إلى العدم؛}{جعل الحب هذا الشعور لطيفًا وخيّرًا، وحوّله الكاهن كله إلى شر.}{}

\Couplet{١٩٠}{فيما يأمرنا العقل بصرامة أن نموت، يتوق الحب إلى حياة وراء القبر؛}{وقلوبنا وعواطفنا وآمالنا ومخاوفنا ستظل تطلب الحياة الآتية.}{}

\Couplet{١٩١}{ومن هنا جاء حلم المستبد الأثير: كنيسة تحكم الدولة وتسوسها؛}{ومن هنا نشأ موكب الأحزان التي لا تُحصى، الملازمة لسلطان الكهنة وحكمهم.}{}

\Couplet{١٩٢}{ومن يجرؤ أن يجيب عن حياة مقبلة؟ ليس لدينا شاهد في مجلس القضاء؛}{إلا ما تخبرنا به كسرة الفخار الشقيقة، من حكايات قديمة وشعوذة جديدة.}{}

\Couplet{١٩٣}{من عاد يومًا ليعلّم الحقيقة، ويرسم لنا أحوال الجنة والجحيم؛}{وكل ما نسمعه لا يصلح إلا لحديث جدّة ونشيد حضانة.}{}

\Couplet{١٩٤}{«ارحمنا، أيها الإنسان!» يصيح الزاهد: «لا تسلبنا أجمل رؤانا؛}{لا بد للبشر من حياة مقبلة توازن أنصبة الحياة غير المتساوية».}{}

\Couplet{١٩٥}{«كلا»، يقول المجوسي، «ليس الأمر كذلك؛ أصب خمري للجميع؛}{كأسًا لهذا وعشرين لذاك، بحسب كِبر مكياله أو صغره».}{}

\Couplet{١٩٦}{«من يشرب كأسًا واحدة فلذته قليلة؛ فقد ولد لأفقر العواطف؛}{ومن يفرغ العشرين عليه دائمًا أن يتوقع الأسف على صداع الصباح».}{}

\Couplet{١٩٧}{يسير آمنًا على الطريق الذي يسمّيه الحكماء الوسط الذهبي؛}{ومن يتسلق جبين جبل أَلْبيّ عابس، لا بد أن يواجه زلات وسقطات كثيرة.}{}

\Couplet{١٩٨}{هنا تلتقي الأطراف؛ فالملوك الممسوحون، الذين تضطرب رؤوسهم المتوّجة؛}{لا تحتوي كأس فرحهم أكثر مما تحتوي كأس المتشردين الذين يموتون على المزابل.}{}

\Couplet{١٩٩}{لخاطئ حكم عليه القدر، ولد ونما ليُعلَّق من خشبة المشنقة؛}{ولقديس يمضي أيامه المقدسة في رجاء منتشٍ بأن يرى إلهه؛}{}

\Couplet{٢٠٠}{ولكل من يتنفس هواءنا الأعلى، توزّع يد القدر دائمًا؛}{وبأجزاء ثابتة متساوية، أنصبتهم من الفرح والحزن، والشقاء والهناء.}{}

\Couplet{٢٠١}{«كيف إذن نقضي فسحة أيامنا في مطاردة الثروة والصيت؛}{ولماذا نجهد، ويجهد البشر جميعًا، من أجل غاية عبثية خيالية؟»}{}

\Couplet{٢٠٢}{الجواب: يطيع البشر ناموسًا يأمرهم بالعمل والكفاح والكد؛}{الحكيم يعرف جيدًا انعدام قيمته، والأحمق يحلم بمكسب أحمق.}{}

\Couplet{٢٠٣}{ومن، حتى بين الحمقى، لا يشعر أن نصف الفرح في السباق}{إلى الثروة والصيت والمكانة، ولا يتنهد حين يأتي النجاح فيتوّج المطاردة؟}{}

\Couplet{٢٠٤}{ومرة أخرى: في الهند أو الصين أو بلاد الفرنج حدثت مصادفة الميلاد؛}{بغير اختيارنا وإرادتنا وكلمتنا؛ والإيمان أيضًا مصادفة.}{}

\Couplet{٢٠٥}{ماذا يقول الفرنجي للهندي: «يا منكِر الشرائع الإلهية؛}{مهما كانت حياتك صالحة تقية، فالجحيم موطنك وموطن أهلك».}{}

\Couplet{٢٠٦}{«اذهب وصفِّ الشراب قبل شربه، واعلم أنك مع كل نفَس تتنفسه؛}{ومع كل خطوة وكل حركة، تميت شيئًا من الحياة».}{}

\Couplet{٢٠٧}{ويجيب الهندي: «امضِ في طريقك الجدير بالملتشا الأدناس الحمقى؛}{واطلب فردوس المنبوذين ذاك وانله؛ أما أنا فأضحك من جنة الكلاب هذه وأبصق».}{}

\Couplet{٢٠٨}{«يا آكلي لحم البقرة المقدسة! يا من تجعلون أفواهكم النهمة قبورًا}{لكائنات لها الحق نفسه في الحياة؛ أي شيطان أباح لكم هذا الإذن القذر؟»}{}

\Couplet{٢٠٩}{وبماذا يصيح الفرنجي في المسلم؟ «أنت موحّد متعدد الزوجات؛}{أعرض عن نبي دجّال، وأحنِ رأسك العنيد ليسوع».}{}

\Couplet{٢١٠}{ويجيب المسلم: «الله واحد، وإن تزوجت أربع مسلمات؛}{أما أنتم فلكم زوجة واحدة، ويا لكم من قوم ملعونين، تقسّمون إلهكم إلى ثلاثة وخمسة».}{}

\Couplet{٢١١}{ويقول البوذي للكونفوشيوسيين: «كالكلاب تعيشون، وكالكلاب تموتون؛}{راضين تستقرون في الأرض البائسة، وتودون تحدّي الله والحساب والجحيم».}{}

\Couplet{٢١٢}{ويرد التتري: «أأقرض نقدي الحاضر الوحيد، الآن؛}{مقابل آنذاك الربوي العبثي الذي عندك؟ ابتعد، يا أحمق ثلاث مرات!»}{}

\Couplet{٢١٣}{«بهذه الحياة الفقيرة، وبهذا العالم الوضيع، أود أن أتمّ ما يكمن فيّ؛}{أسعى إلى إكمال أناي هذه؛ وطموحي الوحيد أن أكون حكيمًا».}{}

\Couplet{٢١٤}{حين يختلف العلماء، من يفصل بينهم وسط الجمع ذي مليارات الرؤوس؛}{ومن سوى المجنون يجرؤ أن يصيح: أنا على صواب، وأنتم جميعًا مخطئون؟}{}

\Couplet{٢١٥}{«أنتم جميعًا على صواب، وأنتم جميعًا مخطئون»، نسمع الصوفي غير المبالي يقول؛}{«فكل واحد يعتقد أن مصباحه المرتعش هو ضوء النهار البهي».}{}

\Couplet{٢١٦}{«لماذا يكون إيمانك باطلًا وإيماني حقًا؟ إنه كله عمل ما لك وما لي؛}{ذلك الحب المولع الأحمق للنفس، الذي يجعل ما لي أفضل مما لك».}{}

\Couplet{٢١٧}{فكفّوا عن مضغ العظام النخرة؛ واسعوا إلى إكساء الهيكل العظمي باللحم والدم؛}{وإلى صوغ صورة يرحب بها الجميع بوصفها جميلة وخيّرة.}{}

\Couplet{٢١٨}{ويقول عربي: «ليس للفتى الكريم من مقام مؤنس إلا الجحيم؛}{أعطونا النار، ولا تعطونا عار صحبة هؤلاء المباركين الحزانى البؤساء».}{\BurtonNote{حاشية برتون: جحيم، أي جهنم أو \textenglish{Gehenna}، النار.}}

\Couplet{٢١٩}{وإذا كانت جنتكم وناركم حقًا، وإذا أجبرني القدر الذي أجبرني على الميلاد}{على الجنة أو النار، فأنا ذاهب؛ وأقابل وقاحة القدر بالازدراء.}{}

\Couplet{٢٢٠}{لا أريد هذا ولا ذاك؛ لقد سئمت من أنا وأنت؛}{وإن تحولنا وتغيرنا كلانا، فماذا يصير من أنت وأنا؟}{}

\Couplet{٢٢١}{يكفي أن نفكر أن أشياء كهذه قد تكون؛ أما القول إنها ليست موجودة أو إنها موجودة؛}{فحمق؛ دعوها جميعًا للقدر، ولا تشنوا على الظلال حربًا عبثية.}{}

\Couplet{٢٢٢}{افعل ما تأمرك به رجولتك، ولا تنتظر استحسانًا إلا من نفسك؛}{أنبل الناس حياةً وأنبلهم موتًا من يضع لنفسه قوانينه ويحافظ عليها.}{}

\Couplet{٢٢٣}{كل حياة أخرى موت حي، عالم لا يسكنه إلا الأشباح؛}{نفَس وريح وصوت ونداء، ورنين جرس الجمل.}{}
